\begin{proposition}[5.8.5]
\label{IV.2.s5.8.5-fr}
Para que un preesquema localmente noetheriano $X$ sea reducido, es necesario
y suficiente que verifique las propiedades $(S_1)$ y $(R_0)$.
\end{proposition}

Teniendo en cuenta (5.7.5), esto resulta de la observación precedente y de
(3.2.1).

\begin{theorem}[5.8.6]
\EGAReaderTarget{ega4.c04.thm.5-8-6}{IV.2.s5.8.6-fr}
(criterio de Serre). Sea $X$ un preesquema localmente noetheriano. Para que
$X$ sea normal, es necesario y suficiente que $X$ verifique las propiedades
$(S_2)$ y $(R_1)$, dicho de otro modo, que para todo $x\in X$ se cumplan las
propiedades siguientes:
\begin{enumerate}
\item[\textup{(i)}] Si $\dim(\mathcal O_x)\leq1$, $\mathcal O_x$ es regular
\textup{(es decir, es un} cuerpo \textup{o un} anillo de valoración discreta \textup{(\textbf{0}, 17.1.4))}.
\item[\textup{(ii)}] Si $\dim(\mathcal O_x)\geq2$, entonces
$\operatorname{prof}(\mathcal O_x)\geq2$.
\end{enumerate}
\end{theorem}

Las condiciones son \emph{necesarias}. En efecto, decir que $X$ es normal significa
que, para todo $x\in X$, $\mathcal O_x$ es un anillo local noetheriano
integralmente cerrado. Si $\dim(\mathcal O_x)=0$ (resp.
$\dim(\mathcal O_x)=1$), se concluye que $\mathcal O_x$ es un cuerpo, puesto
que $\mathcal O_x$ es íntegro (resp. que $\mathcal O_x$ es un anillo de
valoración discreta, en virtud de (\textbf{II}, 7.1.6)). Por otra parte, para todo
elemento $f_x\neq0$ de $\mathcal O_x$, se sabe (Bourbaki, \emph{Alg. comm.},
chap. VII, § 1, n\textsuperscript{o} 4, prop. 8) que los ideales primos asociados a
$\mathcal O_x/f_x\mathcal O_x$ no son inmersos, por tanto $\mathcal O_x$
verifica $(S_2)$ (5.7.7).

Las condiciones son \emph{suficientes}. En efecto, de (5.8.5) resulta primero que
$X$ es \emph{reducido}. Como la cuestión es local, se puede además suponer que
$X=\operatorname{Spec}(A)$, donde $A$ es un anillo noetheriano \emph{reducido}
(\textbf I, 5.1.4); si $R$ es el anillo total de fracciones de $A$, $R$ es producto
directo de un número finito de cuerpos, y (teniendo en cuenta (\textbf{II}, 6.3.6))
bastará probar que $A$ está \emph{integralmente cerrado en} $R$. Sea, pues, $h=f/g$
un elemento de $R$ entero sobre $A$, siendo $g$ y $f$ elementos de $A$ tales
que $g$ no sea divisor de cero. Se tiene una relación de la forma
\begin{equation}
f^n+\sum_{i=1}^n a_i f^{n-i}g^i=0
\qquad\text{con }a_i\in A\quad\text{para }1\leq i\leq n.
\tag{5.8.6.1}\label{IV.2.s5.8.6.1-fr}
\end{equation}

Sea $\mathfrak p$ un ideal primo de $A$ tal que
$\dim(A_{\mathfrak p})=1$; si $f_{\mathfrak p}$ y $g_{\mathfrak p}$ son las
imágenes de $f$ y $g$ en $A_{\mathfrak p}$, de (5.8.6.1) resulta que
$f_{\mathfrak p}/g_{\mathfrak p}$ (que pertenece al anillo total de
fracciones de $A_{\mathfrak p}$, puesto que $g_{\mathfrak p}$ no es divisor
de cero en $A_{\mathfrak p}$ por planitud $(\mathbf{0}_{\mathrm I},5.3.1)$) es entero
sobre $A_{\mathfrak p}$; pero como $A_{\mathfrak p}$ es regular, y por tanto
integralmente cerrado, se tiene
$f_{\mathfrak p}/g_{\mathfrak p}\in A_{\mathfrak p}$. En otras palabras, se
tiene $(fA)_{\mathfrak p}\subset(gA)_{\mathfrak p}$. Pero la hipótesis
$(S_2)$ implica, por (5.7.7), puesto que $g$ no es divisor de cero en $A$, que
$A/gA$ solo tiene ideales primos asociados \emph{no inmersos} $\mathfrak p_i$
($1\leq i\leq n$); ahora bien, $gA$ es la intersección de ideales primarios
$\mathfrak q_i$ correspondientes a los $\mathfrak p_i$, y, según lo que
se acaba de ver, los $\mathfrak q_i$ son las imágenes inversas en $A$, por los
homomorfismos $A\to A_{\mathfrak p_i}$, de los ideales
$(gA)_{\mathfrak p_i}$ (Bourbaki, \emph{Alg. comm.}, chap. IV, § 2, n\textsuperscript{o} 3,
prop. 5). Pero en virtud del Hauptidealsatz (\textbf 0, 16.3.2) se tiene
$\dim(A_{\mathfrak p_i})=1$ para $1\leq i\leq n$, por tanto
$(fA)_{\mathfrak p_i}\subset(gA)_{\mathfrak p_i}$ para todo $i$ según lo
anterior; como $fA$ está contenido en la intersección de las imágenes inversas
de los $(fA)_{\mathfrak p_i}$ ($1\leq i\leq n$), se tiene $fA\subset gA$, es
decir, $f/g\in A$. C.Q.F.D.

\oldpage[IV-2]{109}

\subsection{\texorpdfstring{Módulos $Z$-puros y $Z$-cerrados.}{Módulos Z-puros y Z-cerrados}}

Una parte de las nociones y resultados de esta sección y de la siguiente son
casos particulares de nociones y resultados desarrollados en el capítulo III
en la teoría de la cohomología local. Para comodidad del lector, damos aquí
una exposición independiente.
